\section*{ID10063: Modified Green’s Function: A\_ℓ m(μ\_r)}
\paragraph{Metadata.}
PiID: 5579829368043 | RTEID: RTE:P39076.8638:T1.5548 | UUID: 83646ad7-0c5d-5550-a576-12cfe047eed1

\paragraph{Canonical Statement.}
\textbackslash{}begin\{definition\}[Green’s Function on \textbackslash{}(S\textasciicircum{}2\textbackslash{}times \textbackslash{}mathbb\{R\}\textasciicircum{}+\textbackslash{})] For points \textbackslash{}(r,r'\textbackslash{}) in a half–space bounded by a sphere of radius \textbackslash{}(R\textbackslash{}) with permeability \textbackslash{}(\textbackslash{}mu\_r\textbackslash{}), the Green’s function is \textbackslash{}[ G(r,r') = \textbackslash{}frac\{1\}\{4\textbackslash{}pi |r-r'|\} + \textbackslash{}sum\_\{\textbackslash{}ell=0\}\textasciicircum{}\{\textbackslash{}infty\} \textbackslash{}sum\_\{m=-\textbackslash{}ell\}\textasciicircum{}\{\textbackslash{}ell\} A\_\{\textbackslash{}ell m\}(\textbackslash{}mu\_r) Y\_\{\textbackslash{}ell m\}(\textbackslash{}theta,\textbackslash{}varphi)\textbackslash{},\textbackslash{}overline\{Y\_\{\textbackslash{}ell m\}(\textbackslash{}theta',\textbackslash{}varphi')\} \textbackslash{}frac\{(\textbackslash{}min(r,r')/\textbackslash{}max(r,r'))\textasciicircum{}\{\textbackslash{}ell\}\}\{\textbackslash{}max(r,r')\}. \textbackslash{}] The coefficient \textbackslash{}(A\_\{\textbackslash{}ell m\}(\textbackslash{}mu\_r)\textbackslash{}) is determined by boundary matching and, for homogeneous material parameters, reduces to \textbackslash{}( A\_\{\textbackslash{}ell m\}(\textbackslash{}mu\_r) = \textbackslash{}frac\{(\textbackslash{}ell+1)(\textbackslash{}mu\_r-1)\}\{(\textbackslash{}ell+1)\textbackslash{}mu\_r + \textbackslash{}ell\} R\textasciicircum{}\{2\textbackslash{}ell+1\}. \textbackslash{}) \textbackslash{}end\{definition\}

\paragraph{Canonical Equation.}
\[
A_{\ell m}(\mu_r)
\]

\paragraph{Dictionary.}
Modified Green’s Function: A\_ℓ m(μ\_r) — A\_ℓ m(μ\_r)

\paragraph{Encyclopedia.}
Modified Green’s Function: A\_ℓ m(μ\_r) is a symbolic expression, written A\_ℓ m(μ\_r). It introduces a symbol or relation used elsewhere in the framework. It is built from ℓ, A\_ m, r. Within the Trawin Zero Point registry it serves as one element of the framework's mathematical narrative.
